\begin{center}
\LARGE
Graphic toposes, n-categories and resulting problems of
                       computer algebra
\end{center}

\begin{center}
\Large
F. William Lawvere

\end{center}

\noindent
Date:  July 18th (Thursday\\
Time:  16:45-15:30\\

\begin{center}
\large
Abstract
\end{center}

In 1989 I proposed pre-sheaves of a special kind as algebraic, displayable descriptions of
hierarchical systems. The multi- dimensional graphs underlying n-categories (i.e.ball-and-
hemisphere complexes) form a central example, but finite monoids satisfying the identity xyx = xy
play the pivotal role. New algebraic results involving these toposes include an intimate connection
with Coxeter groups and a characterization of the commonly-occuring case in which the lattice of
sub-toposes is a total order. Computer graphics should be capable of displaying the geometric
realizations of these objects as computer algebra should be capable of solving the relevant word
problems for a known graphic monoid. Whether more than a few such monoids are needed to
govern the modeling of the access-algebra of HYPERTEXT or similar hierarchies remains to be
seen. 

